% 2. Related Work (transferred from docs/paper/paper_draft.md)
\section{Related Work}
\label{sec:02_related_work}

\textbf{Connectome-constrained and bio-inspired network topologies.} Networks whose layer connectivity is defined by biological wiring diagrams have been proposed for the \emph{C. elegans} connectome and compared against classical random graph families (Erdős–Rényi, Watts–Strogatz, Barabási–Albert) \cite{ncnn} \citeneeded{NCNN publication; ER, WS and BA original papers}. Related work on randomly wired architectures suggests that the generator of the wiring pattern, rather than hand design, can drive performance \citeneeded{randomly wired neural networks for image recognition}. That line of work also implies that biological topologies need strong random-graph controls.

\textbf{Drosophila visual connectome models.} The \emph{Drosophila} optic lobe is organised into repeating retinotopic columns containing stereotyped cell types. Its motion-detection pathways (lamina L1–L3; medulla Mi1, Mi4, Mi9, C3; transmedullary Tm1, Tm2, Tm4, Tm9; and the T4/T5 direction-selective neurons) have been characterised from electron microscopy \citeneeded{Takemura et al.; Shinomiya et al.; review of fly motion vision}. Connectome-constrained models of this system have been trained to predict neural activity \cite{flyvis} \citeneeded{FlyVis publication}. Whole-CNS connectomes such as MaleCNS \cite{malecns} and the neuPrint platform \cite{neuprint} now make it possible to extract such circuits programmatically.

\textbf{Graph-structured inductive biases.} Graph neural networks and sparse, graph-masked layers impose relational structure on learned representations \citeneeded{GNN survey or reference}. Whether a \emph{specific} graph helps, beyond its sparsity and degree statistics, is an empirical question that requires matched null graphs.

\textbf{Brain tumour MRI classification.} The Figshare dataset of Cheng et al. \cite{cheng2015enhanced}, \cite{cheng2024dataset} has been widely used for three-class tumour classification (meningioma, glioma, pituitary) \cite{cheng2016retrieval} \citeneeded{representative CNN studies on this dataset}. Many public MRI tumour collections lack patient identifiers, which risks slice-level leakage \citeneeded{data leakage in medical imaging}. The Cheng dataset provides patient IDs and official cross-validation indices \cite{cheng2024dataset}, which makes patient-level evaluation possible.
